\section{Conclusiones}

Esta investigación tuvo el objetivo de cuantificar la reacción competitiva de EDS incumbentes frente a la presencia de nuevos competidores en el mercado chileno observando el comportamiento en precios de sus distintos productos. Bajo metodologías estándar en la literatura relacionada y aproximaciones del márgen de venta se pudo cuantificar la magnitud de la respuesta y caracterízarla. Se encontró que la entrada de un competidor dentro de 2 kilómetros reduce el margen de venta alrededor del $5\%$ en la gasolina 93 y diesel, con efectos no identificables y/o nulos en la gasolina 97. Estimaciones complementarias soportan ideas clave en la literatrura, tales como la atenuación del efecto dada la distancia y  que la respuesta competitiva depende de los márgenes ex-ante al condicionar la respuesta factible. 

El marco teórico que motivó las proposiciones se mantuvo firme frente a las pruebas empíricas. El modelo de \textcite{SALOP1979}, empleado en su formulación original y sin adaptaciones, funcionó bien para describir un mercado cuyo bien es homogéneo, precios altamente visibles y competencia espacial. Las tres hipótesis contrastadas se desprenden de sus resultados de equilibrio, y las tres encuentran respaldo en los datos.

Los resultados se conectan de manera directa con tres definiciones que la Fiscalía Nacional Económica emplea en su labor de fiscalización de este mercado. En primer lugar, la delimitación del mercado relevante geográfico, la \textcite{FNE2021informe} ha trabajado en la práctica con radios de tres a cinco kilómetros. El criterio que la autoridad viene aplicando por analogía y jurisprudencia tiene, entonces, respaldo empírico, y el extremo inferior del rango de tres kilómetros resulta el más apropiado para el análisis de operaciones de concentración.

En segundo lugar, el costo de las barreras a la entrada. La misma Fiscalía documenta que la escasez de terrenos aptos, derivada de las exigencias de los planos reguladores en las zonas urbanas más saturadas, constituye la principal barrera a la entrada de este mercado. Cada entrada que no ocurre son aproximadamente cuatro puntos porcentuales de margen que las incumbentes de ese mercado local retienen de manera indefinida. El resultado sugiere que la regulación de uso de suelo, tiene efectos de competencia mensurables sobre los precios que enfrentan los consumidores.

En tercer lugar, y de manera más específica, la heterogeneidad orienta dónde concentrar el escrutinio. Si la respuesta competitiva se concentra en las estaciones que venían capturando márgenes altos en relación con su entorno, la posición relativa del margen es un indicador observable y de bajo costo para identificar mercados locales donde la entrada tendría mayor efecto disciplinador.

Tres limitaciones conviene reconocer. La primera es que no se dispone de una medida de demanda local a la escala del mercado relevante que permita descartar que la entrada coincida con un aumento de la demanda del sector. La segunda es que el MEPCO es una referencia nacional y no regional, de modo que cualquier diferencial de costo geográfico queda dentro del margen medido. La tercera es que no se observan cantidades vendidas por estación, lo que impide estimar el efecto sobre ingresos y, con ello, evaluar el canal de salida que documentan \textcite{ARCIDIACONO2020} y \textcite{KOH2022}.




